\chapter{Discussion}\label{ch:discussion}

{\doublespacing\normalsize\upshape
\noindent Publication note: Sections~\ref{sec:disc_effects}--\ref{sec:disc_bio} and the associated figures and tables appeared in \citet{bolen_2026} (Bolen, Elzein, Pang, and Hunt, \emph{Actuators}, 2026). Section~\ref{sec:disc_error} is also the subject of an unpublished manuscript in preparation by Bolen, Burgess, Morrow, Elzein, Pang, Scharzenberger, and Hunt \citep{bolen_isometric_2025}. This manuscript has not been submitted for publication. Author contributions for both works are stated in Chapter~\ref{ch:methods}.\par}


\section{Effects of Resting Length}\label{sec:disc_effects}
We have developed new $F_{620}$ equations for \qtylist{10;20}{\mm} diameter Festo BPAs that capture the change in maximum force produced by the actuators at  \qty{620}{\kPa} as a function of their resting length. When examining Equation~(\ref{eq:Fmax10}) and measured data in Figure~\ref{fig:Fmax_fit}b, it can be seen that as $l_0$ goes to infinity, $F_{620_{10}}$ goes to \qty{470}{\N}. When the Festo tool \citep{lang_robert_musclesim_2005} was queried at $l_0=\qty{1}{m}$ as an approximation of infinity, it predicts the $\phi$\qty{10}{\mm} BPA can produce $F=\qty{490}{\N}$ at $P_{620}$. This is within the 10\% variability that Festo states may occur in manufacturing tolerances. Similarly, the Festo data predict a $\phi$\qty{20}{\mm} BPA will produce $F=\qty{1570}{\N}$. Our results for Equation~(\ref{eq:Fmax20}) show that  $F_{620_{20}}$ goes to \qty{1460}{\N} as $l_0$ approaches infinity. Although the difference is much larger, it is also with the \mbox{10\% manufacturing tolerance. }

Where our measured data differs significantly from the Festo model, is that as $l_0$ goes to zero. For the measured data, as the length gets smaller, the force does as well, whereas the Festo tool predicts exponential force increase (see Figure~\ref{fig:Fmax_fit}). As a specific example, when specifying $l_0=\qty{112}{\mm}$ for the \qtylist{10}{\mm} diameter BPAs, the Festo predicted $F_{620}$ is \qtylist{498.6}{\N}. However, using Equation~(\ref{eq:Fmax10}), $F_{620}$ is calculated to be \qtylist{335.3}{\N}. Actual force in the $l_0=\qty{112}{\mm}$ was measured at \qtylist{350.9}{\N} Therefore, the error in predicted $F_{620}$ for the $\phi$\qty{10}{\mm} BPA is $4.5\%$ for our model and $42.1\%$ using the Festo tool. Researchers using BPA resting lengths under \qty{300}{\mm} should take note of these results.

\citet{festo_2026} (p.~12) states that force will be reduced by approximately 10\% for a minimum nominal length. In \citet{festo_2018} (p.~14) they say the force can deviate by 10\% based on factors such as: production variation, material variation, and deviation from nominal length ($l_0$). \citet{festo_2013} (p.~13) notes that testing was done with BPAs that are ten times the uninflated diameter, and that theoretical force can increase by up to 10\%. In each datasheet, the characteristic curves appear to be the same. As demonstrated above, the \qty{100}{\mm} $\phi$\qty{10}{\mm} BPA should have much less force than what is shown in the curve. This is our experience and the reason for the testing. It is unclear why the MuscleSIM tool \citep{lang_robert_musclesim_2005} then predicts higher force as resting length decreases. It could be a reversed sign coefficient in their internal curve-fit model. We have shown here that as the nominal length goes to zero, so does force. 

The reason for BPA models predicting higher force than what is measured experimentally is a recurrent theme across many models: they do not account for the end effects of the pressure vessel where end effects dominate.	Nominally, the clamped ends fix the inner and outer diameter while allowing axial and radial displacement. In the isometric case, there is only radial displacement. The boundary conditions force steep gradients in stress and strain near the ends. Near the boundary conditions are where axial, radial, and circumferential strains deviate from analytic solutions to a cylindrical model. For a constant pressure, the end effect region appears to have the same length. Therefore, in shorter BPAs, the end effects dominate. The axial length scale of the boundary section is comparable to $l_0$,  and the active force differs substantially from that of an assumed perfect cylinder.

Manufacturing variations seem to play an important factor in the behavior of these actuators. It remains to be seen how $\epsilon_{620}$ can be known a priori. We still suspect that it is a function of $l_0$, but it might also be a function of other factors such as product batch number or UV light exposure \citep{festo_2026} (p.~15); \citep{festo_2018} (p.~6); \citep{festo_2013} (p.~13). We found there to be larger variations in the \qtylist{10}{\mm} BPAs than in the \qtylist{20}{\mm} BPAs. Modeling work of ideal McKibben actuators describe how the relationship between contraction and initial length is a function of the total number of twists and the angle of twist in the usable muscle area \citep{chou_measurement_1996}. It is unclear if there are variations in the angle and number of twists in the Festo BPAs due to manufacturing. Uncovering the relationship between $\epsilon_{620}$ and $l_0$ for the Festo DMSP artificial muscles, if it exists in a meaningful way, will require additional controlled tests. 

Similarly, it is not clear how the maximum force might change with manufacturing variations. This was reasonably corrected for by normalizing the equations to each individual's maximum contraction, but more nuanced details might emerge with more data collection. Not knowing these relationships a priori means there is still a need for robots and other systems to go through an iterative design stage if it is discovered that the system does not produce the torque or have the RoM that the design team expects. However, the data and models provided here should enable for faster redesigns with fewer \mbox{additional measurements.}


\section{Nondimensionalized Force Prediction}

We have also introduced the concept of nondimensionalized isometric force that is a function of relative strain and relative pressure (Equation~(\ref{eq:Fstar})). This elegant equation has only three coefficients with low error (see Table~\ref{tab:Fstar}). In the work presented here, we introduce the relative pressure term, $P^*=P/P_{620}$. Our data show that by normalizing force, contraction, and pressure, we are able to create a simplified force equation as a function of contraction and pressure that scales well with initial actuator length, as expressed by Equation~(\ref{eq:force}) in Section~\ref{sec:force_model}.

For example, given $l_0=\qty{54}{\mm}$ and $P=\qty{300}{\kPa}$, the measured force was \mbox{$F_{actual} = \qty{103}{\N}$.} Force prediction using only the Festo tool predicts the force to be \mbox{$F_{predict} = \qty{270}{\N}$.} This is $161\%$ greater than the measured force. Using Equations~(\ref{eq:Fmax10}) and~(\ref{eq:Fstar}) with Equation~(\ref{eq:force}) gives $F_{predict} = F^*(0, 0.48) \cdot F_{620}(0.054) = 0.49 \cdot \qty{230}{\N} = \qty{112}{\N}$. This is an error of $8\%$, which is much more accurate than using the Festo tool only.

The presented work is capable of predicting force output for 10 mm and 20 mm muscles with different initial muscle lengths, internal pressure, and contraction amount with minimal initial data collection. However, it remains to be seen if these relationships are maintained over the full lifecycle of the muscles. There are not yet any peer reviewed studies investigating the long-term behavior of these artificial muscles, though the Festo datasheet provides some indication of what can be expected \citep{festo_2026} (p.19). They report that the 40 mm BPA can complete 1 million cycles at \qty{6000}{\N} (maximum rated load), or 10 million cycles at \qty{4000}{\N}. The service life of the other size BPAs is not reported; however, it is mentioned that thermal stressing may reduce the lifespan and they should be inspected every 500,000 cycles for signs of aging \citep{festo_2018} (pp.~10--11). Future research should look at the behavior of these muscles over time.

In our work, $P_{620}=\qty{620}{\kPa}$ is the maximum supply pressure for our system, and other users of this actuator may use a different supply pressure. However, this supply pressure is not required, and normalizing by this value was chosen as a method for improving the optimization techniques by reducing the sensitivity of the term associated with pressure. This equation will work for any pressure supplies below \qty{620}{\kPa}. We also anticipate the equation will also work for pressures up to \qty{800}{\kPa} (the maximum rated pressure for Festo $\phi$\qty{10}{\mm} BPAs); however, we do not have the data to confirm this.

\section{Comparison with Other Published BPA Models}

The proposed model was compared with the formulations by \citet{sarosi_new_2012}, \citet{sarosi_10mm}, and \citet{martens_modeling_2017}, to evaluate their predictive performance across actuator sizes and operating conditions (see Table~\ref{tab:model_compare} and Figure~\ref{fig:ModelCompare}).
For the $\phi$\qty{20}{\mm} BPA, parameters from \cite{sarosi_new_2012} were used, corresponding to actuators with resting lengths of \qty{200}{\mm} and \qty{400}{\mm}. These were compared against the experimental data at $l_0=\qty{300}{\mm}$ and $l_0=\qty{450}{\mm}$, respectively. For the $\phi$\qty{10}{\mm} BPA with resting lengths of \qty{257}{\mm} and \qty{233}{\mm}, parameters from \cite{sarosi_10mm} were applied. In each case, the maximum contraction parameter ($k_{max}$) was determined numerically by solving $F(\qty{6.2}{\bar}, k_{max}) = 0$, ensuring consistency with the experimental maximum pressure condition.




\begin{table}[H]
\caption{Goodness-of-fit comparison for the Bolen, Sárosi, and Martens and Boblan models at different resting lengths (in mm) for \qty{10}{\mm} and \qty{20}{\mm} BPAs. RMSE and maximum error are reported in newtons, and FVU denotes the fraction of variance unexplained.}
\label{tab:model_compare}
\tiny
\begin{adjustwidth}{-2cm}{0cm}
%\centering %% If there is a figure in wide page, please release command \centering
\begin{tabularx}{\textwidth}{@{}lllllllllll@{}}
\toprule
\textbf{Diameter} &
  \textbf{Length} &
  \textbf{} &
  \textbf{Bolen} &
  \textbf{} &
  \textbf{} &
  \textbf{Sarosi} &
  \textbf{} &
  \textbf{} &
  \multicolumn{2}{l}{\textbf{Martens and Boblan}} \\ 
\textbf{} &
  \textbf{} &
  \textbf{RMSE} &
  \textbf{FVU} &
  \textbf{Max. Error} &
  \textbf{RMSE} &
  \textbf{FVU} &
  \textbf{Max Error} &
  \textbf{RMSE} &
  \textbf{FVU} &
  \textbf{Max. Error} \\ \midrule
\textbf{10 mm} & 257 & 22.9 & 0.0353 & 46.0  & 143.8 & 1.3903 & 235.8 & 73.8  & 0.3668 & 122.9  \\
\textbf{}      & 233 & 16.8 & 0.0187 & 30.3  & 132.7 & 1.1655 & 222.6 & 62.6  & 0.2590 & 101.2  \\ \midrule
\textbf{20 mm} & 300 & 35.6 & 0.0144 & 114.0 & 132.1 & 0.1992 & 313.4 & 72.8  & 0.0605 & 137.8  \\
\textbf{}      & 450 & 67.6 & 0.0258 & 166.6 & 177.7 & 0.1785 & 294.8 & 862.1 & 4.2022 & 1321.1 \\ \bottomrule
\end{tabularx}
\end{adjustwidth}
\end{table}

While the \citet{sarosi_new_2012} and \citet{sarosi_10mm} models capture general trends in force–contraction behavior, discrepancies were observed when applying parameters derived from different actuator lengths. In particular, the model tends to over-predict force at higher pressures and low amounts of contraction, suggesting sensitivity to geometric scaling effects not explicitly accounted for in the formulation.

The \citet{martens_modeling_2017} model  is derived from a geometric and material-based framework and provides a more physically grounded representation of actuator behavior. However, its accuracy depends strongly on precise geometric parameters, including braid angle, wall thickness, and fiber length. Small deviations in these parameters can lead to significant differences in predicted force, particularly at higher pressures.

\begin{figure}[H]%[tbp]
%\begin{adjustwidth}{-\extralength}{0cm}
%\centering
%\begin{center}
    \includegraphics[width=.99\textwidth]{05_Comparison.eps}
%\end{center}
%\end{adjustwidth}
\caption{Comparison between measured and predicted BPA force at varying amounts of contraction and pressure. A $y=x$ perfect fit line is added for reference. The BPAs compared are: (\textbf{a}) $\phi$\qty{10}{\mm}, $l_0=\qty{0.257}{\mm}$, (\textbf{b}) $\phi$\qty{10}{\mm}, $l_0=\qty{0.233}{\mm}$, (\textbf{c}) $\phi$\qty{20}{\mm}, $l_0=\qty{0.300}{\mm}$, and (\textbf{d}) $\phi$\qty{20}{\mm}, $l_0=\qty{0.450}{\mm}$.}\label{fig:ModelCompare}
\end{figure}

The Martens and Boblan model was first evaluated at the same pressure–contraction points as in their experimental data \citep{martens_modeling_2017} (Tables A1 and  A2), indicating that it reproduces the general published force map shape but is sensitive to the assumed initial geometry. Examining their $\phi$\qty{20}{\mm} BPA, for example, using the explicitly stated $l_0=\qty{300}{\mm}$ and initial diameter $D_{0}=\qty{20}{\mm}$ significantly under-predicted the results compared to using the implied $l_0=\qty{296}{\mm}$ (\citep{martens_modeling_2017} (Table A1)) and $D_{0}+H_{0}=\qty{21.8}{\mm}$, where $H_{0}=\qty{1.8}{\mm}$ is the material thickness. The diameter adjustment can be implied \mbox{from \citep{chou_measurement_1996}.} This adjusted diameter was also included for the $\phi$\qty{10}{\mm} BPA to give a better fit. As can be seen in Figure~\ref{fig:ModelCompare} and Table~\ref{tab:model_compare}, the highest accuracy is observed when $\phi$\qty{20}{\mm} and $l_0=\qty{300}{\mm}$, which is at or very near their experimental measurements. For the $\phi$\qty{10}{\mm} BPAs the accuracy decreases even though their experimentally measured $l_0=\qty{250}{\mm}$ is close to our $l_0=\qty{257}{\mm}$ and $l_0=\qty{233}{\mm}$. The largest error occurs at $\phi$\qty{20}{\mm} and $l_0=\qty{450}{\mm}$ and the model starts to predict negative force. We have calculated that the $F_{PE}\cdot\frac{dD}{DL}$ term becomes very large compared to the others \citep{martens_modeling_2017} (p.~5). This suggests that the model is sensitive to geometric parameter selection and interpretation, particularly with respect to initial length and diameter definitions.

In comparison, the present normalized model demonstrates improved consistency across actuator lengths by incorporating experimentally derived scaling through $F_{620}$ and $\epsilon_{620}$. This normalization reduces sensitivity to geometric variation and enables a unified representation of actuator behavior, while maintaining predictive accuracy over the tested range of pressures and contractions.


\section{Implications for Real-Time Control and Simulation}

Although pneumatic muscle actuators exhibit hysteresis, thermodynamic effects, and other dynamic behavior, accurate and high-performance controller designs can still be achieved using a static force map approach when the underlying static force model is accurate. Martens and Boblan explicitly note that BPA dynamics are often dominated by the air dynamics inside the actuator and the mass flow behavior of the valve, such that precise modeling of the static force characteristic remains crucial for torque control \mbox{accuracy \citep{martens_modeling_2017}.} This view is consistent with more recent BPA control work, including a sensor-less torque control interface for BPA-driven joints \citep{martens_framework_2018}, stable force/position/stiffness control of antagonistic pneumatic-muscle joints \citep{ugurlu2019stable, hunt_development_2017}, and sensor-less angle stiffness control of antagonistic BPA systems \citep{shinSensorlessAngleStiffness2022}. 

In this context, the present model is well suited for real-time use because it reduces the static isometric force prediction to a low-order normalized surface, $F^*(\epsilon^*,P^*)$, scaled by the resting-length-dependent term $F_{620}(l_0)$. This structure is computationally inexpensive to evaluate online and is also suitable for lookup-table implementation. While the present work does not model the full actuator--valve dynamics, it provides the static force relationship needed as a feedforward or inner-model component in real-time torque control, simulation, and preliminary design workflows. It is anticipated that improved real-time control for high-speed and high-precision systems will benefit more from improved modeling of airflow dynamics than dynamics involving actuator hysteresis, damping, and other dynamic behavior. Characterizing the actuators under dynamic operating conditions --- isokinetic (constant-velocity), isobaric, and isotonic (constant-load) contractions --- is a natural next step beyond the isometric characterization reported here, for which the pressure--force dynamics model of \citet{elzein_pressure_2025} provides the modeling foundation.



\section{Biomimetic Artificial Muscle Analogy}\label{sec:disc_bio}

BPA artificial muscles are often said to be analogous to biological muscles because they have force--length curves and can produce force only in tension. This analogy is worth closer inspection. For example, the describes the passive force–length curve has been described as an exponential function arising from connective tissue elements, while the active force–length relationship is typically modeled as a bell-shaped curve reflecting actin–myosin overlap \citep{winters1990}. These formulations emphasize that muscle force generation is governed by nonlinear geometric and material interactions rather than simple linear elasticity.
McKibben-type pneumatic artificial muscles (of which Festo is an example) exhibit analogous nonlinear behavior, though arising from fundamentally different physical mechanisms. As shown by Chou and Hannaford, the force–length relationship of McKibben actuators is derived from geometric constraints imposed by the braided shell, leading to a nonlinear dependence of force on actuator length through the braid angle \cite{chou_measurement_1996}. In their formulation, actuator force is proportional to pressure and a nonlinear geometric term $(3\cos^2\theta - 1)$, which governs both the maximum contraction limit and the curvature of the force–length relationship. The work presented here demonstrates that this nonlinear curve can be well approximated by an exponential as well.

However, there are some distinct differences between how these artificial muscles behave when compared to biological muscles.
For example, the biological muscle optimum fiber length $l_{OFL}$ is not equivalent to $l_0$ in the muscles. The maximum active force that the artificial muscles can produce is at the resting length, also the maximum length, of the artificial muscles. In contrast, the maximum active force a biological muscle length can produce is somewhere in the middle of the total contractile length of the muscle. 

Furthermore, biological muscles can adapt the optimal operating length through structural remodeling (e.g., sarcomere addition or tendon compliance) while BPAs have a fixed geometric configuration. This allows biological systems to self-optimize for the particular tasks the animal is subjected to. In order to re-optimize a robot for new tasks, new actuators would have to be swapped into the system. This distinction highlights a key limitation of BPAs in biomimetic applications, while also reinforcing the importance of accurate nonlinear modeling for control and design.




\section{Sources of Error in Isometric Torque Generation}\label{sec:disc_error}

Our improved BPA characterization predicts torque more accurately than the original BPA model (see Fig.~\ref{fig:FlxPin_group}\textbf{A} and \ref{fig:FlxPin_group}\textbf{B}). This only accounts for part of the error. There are certainly many factors that affect the torque results. Compliance in the testing system, position location tolerance, flexibility of brackets, imperfect path modeling, and variation in individual BPAs are other factors.

Simplifying the model and testing it allowed us to see how we were deficient in our previous analysis. The isometric system is not rigid. There is compliance in the winch cable, artificial tendons, and brackets that bend and stretch. We introduced a constant length offset term $\chi_{0}$, axial stiffness $\chi_{1}$, and bending stiffness $\chi_{2}$. The bracket frame of reference was placed where it starts to cantilever and is rotated about $Z$.  Optimization found results for a $\phi$\qty{10}{\mm} BPA with $l_{0}=\qty{48}{\cm}$ in the simple pinned knee configuration. These results were validated with second $\phi$\qty{10}{\mm} BPA with $l_{0}=\qty{45.7}{\cm}$.

The modeling of extensor BPA paths that wrap around the knee joint is an approximation with inherent error (Fig. \ref{fig:Ext_passthru}). When the moment arm is short, small changes to it can have a large effect on the resultant torque \citep{young_analyzing_2019}. It was assumed that wrapping the BPA around the joint would be easier to account for than adding a patella, deterministically designing its location as a function of knee angle, and accounting for friction. During greater magnitude knee flexion angles, shorter $l_{0}$ BPAs were noted as being stretched and compressed (see Fig.~\ref{fig:testJigs}\textbf{C}). The same force that caused the stretch and compression would also sometimes cause the BPA to slip off of the anterior bolt head of the bolt that holds the femur to the knee. This would result in a $\pm \; 20\,\text{mm}$ displacement in the $Z$ direction, which means a change the force vector and a change in torque. The system designer, too, must be mindful. Fig. \ref{fig:ExtPin_group} shows a deviation between measured and originally predicted results that widens as expected torque increases. This indicates system compliance. However, the  hybrid calculated torque matching or exceeding predicted torque, except for higher degrees of flexion in the $l_{0}=\qty{42}{\mm}$, indicates a loss of functional $l_m$. BPA This is why the $\chi_{3}$ term was created in \eqref{eq:deltaL}. Fig. \ref{fig:ExtPin_group} shows the improved fit from running the optimization that includes this term.

The biomimetic knee joint with an extensor BPA was used to validate the extensor optimization results. The predicted line more closely matches the torque values. However, the predicted torque is still higher than experimentally measured torque.


In the presented nonlinear biomimetic robotic system, rigid body mechanics simplifications, often encountered in other robotic and biomechanical systems, are removed. Instead, working from the simple to the complex elucidated mechanisms of system compliance that will lead to much more accurate first iteration designs. Technologies have been developed to predict system bending and axial compliance, a constant length offset, and a loss of usable length that occurs as artificial muscle bends around a joint. 

In equation $m\ddot \varepsilon^* + c (\dot \varepsilon^* ) \dot \varepsilon^* + k ( \varepsilon^* ) \varepsilon^* = F$ given $k(\varepsilon^{*}, P^{*})
= \left.\frac{\partial F}{\partial \varepsilon^{*}}\right|_{P^{*}}
= -c_{0}c_{1}e^{-c_{1}\varepsilon^{*}}
  - 2c_{2}\varepsilon^{*}P^{*}e^{-c_{2}(\varepsilon^{*})^{2}}$
then if we set mass m and external force F to zero we get
$c (\dot \varepsilon^* ) = (k ( \varepsilon^* ) \varepsilon^*)/\dot \varepsilon^*$


\section{Synthesis: From Actuator Characterization to Neuromechanical Control}\label{sec:synthesis}

Taken together, the actuator and joint level results of this dissertation form a coherent design chain for biomimetic legged robots. At the actuator level, the maximum isometric force of Festo BPAs was shown to depend strongly on resting length, correcting a systematic over-prediction by the manufacturer's tool at short resting lengths, and a nondimensional force surface was derived that predicts isometric force from relative strain and relative pressure with only three diameter-dependent coefficients. At the joint level, the improved actuator model was combined with a compliance-corrected musculotendon kinematics model to predict isometric knee torque across the full range of motion for pinned and biomimetic knee joints, and the biomimetic knee driven by optimized BPAs was shown to meet or exceed human isometric torque over most of the range of motion.

This chain was developed deliberately so that each stage can feed the next. The force model supplies the actuator terms for the torque model; the torque model supplies the plant model for designing joint controllers; and the validated biomimetic knee demonstrates that the actuator architecture can reproduce human joint-level force capabilities, which is a prerequisite for the neuromechanical control studies planned in Chapter~\ref{ch:futurework}. What remains missing is the neural layer itself: a controller that schedules the measured muscle forces the way the nervous system schedules motor unit activity. The AARL's prior AnimatLab walker (Section~\ref{sec:prior_walker}) offers an early demonstration of both halves of that challenge: acting as a feedforward controller, its two-layer CPG drove gait-like stepping of a suspended biomechanical model --- confirming that rhythm generation does not require sensory feedback --- but untuned afferent feedback could not support those patterns against the ground, which is why feedback tuning is the critical path from feedforward stepping to weight-bearing walking. Closing that loop in simulation, and finally on hardware, is the subject of the next chapter.
